% This Source Code Form is subject to the terms of the Mozilla Public
% License, v. 2.0. If a copy of the MPL was not distributed with this
% file, You can obtain one at https://mozilla.org/MPL/2.0/.

\documentclass[12pt, a4paper]{article}
\usepackage[T2A]{fontenc}
\usepackage[utf8]{inputenc}
\usepackage[russian]{babel}
\usepackage{geometry}
\geometry{left=2cm, right=2cm, top=2cm, bottom=2cm}
\usepackage{array}
\usepackage{multirow}
\usepackage{amsmath}
\usepackage{indentfirst}

\title{Лабораторная работа №1}
\author{Дыкова А.В., группа 5030302/10801}
\date{18.09.2023}

\begin{document}
	
	\subsubsection*{Опыт №1}
	\subsubsection*{Применение индикаторов для определения pH раствора}
	
	Мы рассмотрели 9 колб с различными средами, в которые предварительно были добавлены индикаторы.
	
	\begin{center}
		\begin{tabular}{|c|c|c|c|}
			\hline
			Индикатор & Кислая среда & Нейтральная среда & Щелочная среда \\
			& pH < 7 & pH = 7 & pH > 7 \\
			\hline
			Метилоранж & розовый & оранжевый & жёлтый \\
			\hline
			Лакмус & красный & фиолетовый & синий \\
			\hline
			Фенолфталеин & бесцветный & бесцветный & малиновый \\
			\hline
		\end{tabular}
	\end{center}
	
	\textbf{Вывод:} Наиболее универсальный индикатор -- лакмус, с помощью него можно определить реакцию среды любого раствора (кроме цветных и очень разбавленных). С помощью фенолфталеина можно определить только щелочную реакцию среды. Метилоранж используется в химии как вспомогательный индикатор.
	
	\subsubsection*{Опыт №2}
	\subsubsection*{Получение и свойства амфотерного гидроксида (амфолита)}
	
	Мы внесли в две пробирки по 2 капли раствора гидроксида натрия. Затем мы добавили в каждую по 5 капель раствора хлорида алюминия. При этом образовался белый осадок:
	
	1) $3\ NaOH + AlCl_{3} = Al(OH)_{3} \downarrow + 3\ NaCl$ 
	
	В первую пробирку к полученному осадку мы добавили 10 капель раствора соляной кислоты и наблюдали растворение осадка:
	
	2) $Al(OH)_{3} + 3\ HCl = AlCl_{3} + 3H_{2}O$ 
	
	Во вторую пробирку к полученному осадку мы добавили 10 капель раствора гидроксида натрия и наблюдали растворение осадка:
	
	3) $Al(OH)_{3} + 3\ NaOH = Na[Al(OH)_{4}] + 3H_{2}O$ 
	
	\textbf{Вывод:} Гидроксид алюминия является амфотерным. Его основные свойства проявляются при взаимодействии с кислотами. В результате образуется соль, в которой алюминий является катионом. Кислотные свойства гидроксида алюминия проявляются при взаимодействии с щелочами. В результате образуется соль, в которой алюминий находится в составе аниона.
	
	\subsubsection*{Опыт №3}
	\subsubsection*{Получение кислой соли (гидроксосоли) и перевод её в нейтральную (среднюю).}
	
	Мы налили в большую пробирку раствор гидроксида кальция, а затем пропускали в неё диоксид углерода из аппарата Киппа до выпадения осадка:
	
	1) $Ca(OH)_{2} + CO_{2} = CaCO_{3} \downarrow + H_{2}O$ 
		
	После получения осадка мы продолжили пропускать диоксид углерода из аппарата Киппа до полного растворения осадка:
	
	2) $CaCO_{3} + (CO_{2} + H_{2}O) = Ca(HCO_{3})_{2}$ 
	
	На следующем этапе полученный раствор кислой соли мы разлили в две пробирки поровну.
	
	В первую пробирку мы добавили гидроксид кальция и наблюдали помутнение раствора (выпадение осадка)
	
	3) $Ca(HCO_{3})_{2} + Ca(OH)_{2} = 2CaCO_{3} \downarrow + 2H_{2}O$
	
	Во вторую пробирку мы добавили примерно столько же гидроксида натрия и наблюдали меньшее помутнение раствора:
	
	4) $Ca(HCO_{3})_{2} + 2NaOH = CaCO_{3} \downarrow + Na_{2}CO_{3} + 2H_{2}O$
	
	\textbf{Ответы на вопросы:}
	
	1. Уравнения реакций, описывающих химические процессы в опыте, представлены под номерами 1-4.
	
	2. Химическая сущность превращения кислой соли в нормальную заключается в нейтрализации ионов водорода $H^{+}$, входящих в состав кислотных остатков кислых солей, гидроксид-ионами ${OH}^{-}$ с образованием малодиссоциированных молекул воды.
	
	3. Другие способы получения кислых солей:
	
	а) $2H_{2}SO_{4} + Ba(OH)_{2} = Ba(HSO_{4})_{2} + 2H_{2}O$
	
	б) $2SO_{3} + Ba(OH)_{2} = Ba(HSO_{4})_{2}$
	
	в) $2H_{2}SO_{4} + BaO = Ba(HSO_{4})_{2} + H_{2}O$
	
	\subsubsection*{Опыт №4}
	\subsubsection*{Получение основной соли (гидроксосоли) и перевод её в нормальную (среднюю)}
	
	Мы влили в пробирку 5 капель раствора сульфата меди (II) и добавили такой же объём раствора гидроксида натрия, при этом выпал голубой осадок:
	
	1) $CuSO_{4} + 2NaOH = Cu(OH)_{2} \downarrow + Na_{2}SO_{4}$ 
	
	Затем мы добавили ещё 12 капель сульфата меди при постоянном помешивании палочкой и наблюдали изменение цвета осадка на другой оттенок голубого:
	
	2) $Cu(OH)_{2} + CuSO_{4} = (CuOH)_{2}SO_{4} \downarrow$
	
	К полученному осадку основной соли мы добавили 10 капель серной кислоты и наблюдали полное растворение осадка:
	
	3) $(CuOH)_{2}SO_{4} + H_{2}SO_{4} = 2CuSO_{4} + H_{2}O$ 
	
	\textbf{Ответы на вопросы:}
	
	1. Уравнения реакций, описывающих химические процессы в опыте, представлены под номерами 1-3.
	
	2. Другие способы получения основных солей:
	
	а) $H_{2}SO_{4} + 2Ba(OH)_{2} = (BaOH)_{2}SO_{4} + 2H_{2}O$ 
	
	б) $SO_{3} + 2Ba(OH)_{2} = (BaOH)_{2}SO_{4} + H_{2}O$
	
	в) $H_{2}SO_{4} + 2BaO = (BaOH)_{2}SO_{4}$
	
	3. Химическая сущность превращения основной соли в нормальную заключается в нейтрализации гидроксид-ионов ${OH}^{-}$, входящих в состав катионов основных солей, ионами водорода $H^{+}$ с образованием малодиссоциированных молекул воды.
	
	\subsubsection*{Выводы}
	В результате первой лабораторной работы мы изучили поведение трёх наиболее часто применяемых индикаторов (лакмуса, метилоранжа и фенолфталеина) в различных растворах неорганических соединений (кислых, нейтральных и щелочных). Также, мы проанализировали химические свойства амфотерного гидроксида и узнали, что он проявляет кислотные свойства при взаимодействии с щелочами и щелочные при взаимодействии с кислотами. Кроме того, мы познакомились с методами получения и нейтрализации кислых и основных солей.
	
\end{document}